% =========================================================
% frontmatter/acknowledgements.tex
% Ucapan terima kasih.
% =========================================================

\chapter*{Ucapan Terima Kasih}
\addcontentsline{toc}{chapter}{Ucapan Terima Kasih}
\markboth{Ucapan Terima Kasih}{Ucapan Terima Kasih}

Penyusunan buku ajar ini tidak terlepas dari dukungan berbagai
pihak. Penulis menyampaikan terima kasih kepada:

\begin{enumerate}
    \item \BookProgram\ dan \BookFaculty, \BookInstitution, atas
          dukungan dan fasilitas yang diberikan selama proses
          penyusunan buku ajar ini.

    \item Tim penyusun kurikulum berbasis \textit{Outcome-Based
          Education} program studi, atas kerangka capaian
          pembelajaran dan dokumen Rencana Pembelajaran Semester
          yang menjadi acuan buku ini.

    \item Para dosen dan asisten laboratorium yang telah memberikan
          masukan terhadap materi praktikum, urutan penyajian, dan
          pemilihan kasus pemrograman.

    \item Mahasiswa peserta mata kuliah \BookCourse\ yang telah
          memberikan umpan balik atas kejelasan materi, tingkat
          kesulitan latihan, dan keefektifan contoh program.

    \item Penulis-penulis sumber terbuka yang karyanya menjadi
          rujukan buku ini, sebagaimana tercantum pada Daftar
          Pustaka dan Lampiran~D.
\end{enumerate}

Semoga buku ini bermanfaat bagi mahasiswa dan siapa pun yang ingin
memahami cara kerja komputer pada level instruksi mesin.

\vspace{1cm}

\begin{flushright}
\begin{tabular}{c}
Penulis
\end{tabular}
\end{flushright}

\clearpage
